\scanpage{264}

\deleted{$\lambda\in T_H$ tel que $\unsure{\lambda u(\sigma)|\widetilde\lambda u(\sigma)}U_0=$}

pond à $\sigma.u(\lambda)U_0$, donc la condition

\[gt_0=\lambda t_0\]

s'écrit aussi

\[\bar g(\sigma)U_0=\sigma.u(\lambda)U_0\qquad\begin{aligned}&\text{pour }\lambda\in Z\text{ convenable}\\&\text{(d'ailleurs unique)}\end{aligned}\tag*{(13\unsure{8|6})}\]

En fait, ce\unsure{tte|s} conditions impliquent déjà la condition, en apparence plus forte, que $\bar g$ \inserted{(normalise $U_0$ et)} \uline{centralise} $T_H$ lui-$=$, et que de plus on a

\[\bar g(\sigma)=\sigma u(\lambda)\qquad\text{pour }\lambda\in Z\text{ convenable},\tag*{(13\unsure{9|8})}\]

soit encore

\[\bar g(\sigma)\in\sigma.T_H,\tag*{(1\unsure{40|39})}\]

\deleted{puisque} \inserted{En effet,} en termes de $g\in\inserted{T_0\subset}G\subset\unsure{\mathfrak S_S|\mathscr G_S}$, \deleted{cette relation} \inserted{on sait} bien que $g$ centralise $T_H$, et (139) s'écrit \deleted{s'écrit}

\[g\sigma g^{-1}=\sigma u(\lambda)\qquad\text{i.e.}\qquad\sigma(g)g^{-1}=u(\lambda)\]

ou encore, \deleted{en écrivant $t$ au lieu de $g$ (puisque $g\in T_0$)}

\[\overbrace{\sigma(\deleted{t}g)}^{\rotatebox{270}{$\in$}\, T_1}=u(\lambda)\overset{\substack{T_0\\\rotatebox{270}{$\in$}}}{\deleted{\unsure{t}}g}\qquad \deleted{\text{avec }\lambda=\delta(g)},\ \substack{\text{pour }g\in T_0\\\text{et }\lambda\in Z\text{ conv.}}\tag*{(141)}\]

\deleted{soit encore, posant $\widetilde\delta(t)=\widetilde\lambda$}

\[\deleted{u_1(\widetilde\lambda)=u(\lambda)u_0(\widetilde\lambda)}\]

\deleted{On a vu} \inserted{En effet,} les deux mbes sont dans $T_G$ (qui est normalisé par $\sigma$), et on a vu qu'ils sont congrus mod $U_0$, donc congrus mod $T_G\cap U_0=T_H\cap U_0=\{1\}$, donc ils sont égaux. De plus, \deleted{en prenant appliquant} \inserted{cette formule montre} que \inserted{$g$ et} $\sigma(g)$ \deleted{et $g$} ont $=$ image \inserted{$\widetilde\lambda$} dans $\widetilde Z=T_G/T_H$, donc, \deleted{comme $g$ \unsure{prouve}} \inserted{comme ils sont} dans $T_0$ et $T_1$ respectivement, ils peuvent s'écrire $u_0(\widetilde\lambda),u_1(\widetilde\lambda)$, et (141) s'écrit

\[u_1(\widetilde\lambda)=u(\lambda)u_0(\widetilde\lambda)\tag*{(142)}\]

qui montre qu'on a (cf (119))

\[\lambda=c(\widetilde\lambda).\tag*{(143)}\]

% archive sheet number: 263
